% ============================================================
%  BAB 3 – METODOLOGI PENELITIAN
%  Compile standalone:  pdflatex bab-3.tex
%  Compile as part of thesis: pdflatex main.tex
% ============================================================
\documentclass[main]{subfiles}
\usepackage{hyphenat}
\usepackage{iftex}

% ── Encoding, Language & Fonts ───────────────────────────────
\ifXeTeX
    \usepackage{polyglossia}
    \setmainlanguage{bahasai}
    \usepackage{fontspec}
    \setmainfont{TeX Gyre Termes}
    \setsansfont{TeX Gyre Heros}
    \setmonofont[Scale=0.88]{TeX Gyre Cursor}
\else
    \usepackage[utf8]{inputenc}
    \usepackage[T1]{fontenc}
    \usepackage[indonesian]{babel}
    \usepackage{mathptmx}
    \usepackage{helvet}
    \usepackage{courier}
\fi

% ── Page geometry (A4, UI thesis margins) ────────────────────
\usepackage[
  a4paper,
  top    = 2.54cm,
  bottom = 2.54cm,
  left   = 3.17cm,
  right  = 2.54cm
]{geometry}

% ── Line spacing & paragraph ─────────────────────────────────
\usepackage{setspace}
\onehalfspacing
\setlength{\parindent}{1.27cm}
\setlength{\parskip}{0pt}

% ── Microtype & justification ────────────────────────────────
\usepackage{microtype}
\sloppy

% ── Headings ─────────────────────────────────────────────────
\usepackage{titlesec}

\titleformat{\chapter}[block]
  {\centering\sffamily\bfseries\large}{}
  {0pt}{\MakeUppercase}
\titlespacing*{\chapter}{0pt}{0pt}{1.5em}

\titleformat{\section}
  {\sffamily\bfseries\normalsize}{\thesection}{1em}{}
\titlespacing*{\section}{0pt}{1.2em}{0.6em}

\titleformat{\subsection}
  {\sffamily\bfseries\normalsize}{\thesubsection}{1em}{}
\titlespacing*{\subsection}{0pt}{1em}{0.5em}

\titleformat{\subsubsection}
  {\sffamily\bfseries\normalsize}{\thesubsubsection}{1em}{}
\titlespacing*{\subsubsection}{0pt}{0.8em}{0.4em}

% ── Tables ───────────────────────────────────────────────────
\usepackage{longtable, booktabs, array, multirow}
\usepackage{calc}

% ── Math ─────────────────────────────────────────────────────
\usepackage{amsmath, amssymb}

% ── Graphics ─────────────────────────────────────────────────
\usepackage{graphicx}
\makeatletter
\def\maxwidth{\ifdim\Gin@nat@width>\linewidth\linewidth\else\Gin@nat@width\fi}
\makeatother
\setkeys{Gin}{width=\maxwidth,keepaspectratio}

% ── TikZ (ontology diagram) ──────────────────────────────────
\usepackage{tikz}
\usetikzlibrary{shapes.geometric, arrows.meta, positioning, calc, fit, backgrounds}

% ── Extended tables ──────────────────────────────────────────
\usepackage{tabularx}

% ── Lists ────────────────────────────────────────────────────
\usepackage{enumitem}
\setlist[enumerate,1]{
  label      = \arabic*.,
  leftmargin = 1.27cm,
  itemsep    = 4pt,
  parsep     = 0pt,
  topsep     = 4pt
}
\setlist[enumerate,2]{
  label      = \alph*.,
  leftmargin = 1.27cm,
  itemsep    = 3pt,
  parsep     = 0pt,
  topsep     = 2pt
}

% ── Hyperlinks ───────────────────────────────────────────────
\usepackage[hidelinks]{hyperref}

% Bibliography setup is inherited from main.tex (biblatex + references.bib)
% when this file is compiled standalone via the subfiles class.

% ============================================================
\begin{document}

% When standalone: set counter so \chapter{} renders as "BAB 3".
% When in main.tex: main.tex sets the counter before \subfile{bab-3}.
\ifSubfilesClassLoaded{}{\setcounter{chapter}{2}}

\chapter{Metodologi Penelitian}

Bab ini menjabarkan kerangka kerja yang digunakan untuk membangun dan
menguji sistem \textit{Knowledge Graph Completion} (KGC) berbasis
\textit{Large Language Model} (LLM), mencakup desain penelitian, sumber
dan pengumpulan data, kerangka konsep dan skema ontologi, parameter
sistem, serta prosedur evaluasi terhadap hasil yang dihasilkan.

% ============================================================
\section{Desain Penelitian}
% ============================================================

Penelitian ini menggunakan pendekatan \textit{Design Science Research}
(DSR) yang difokuskan pada pengembangan purwarupa sistem KGC berbasis
LLM. DSR dipilih karena penelitian ini secara spesifik bertujuan
menghasilkan artefak inovatif, representasi pengetahuan kurikulum yang
bersifat lintas-disiplin, sebagai respons terhadap permasalahan
fragmentasi pembelajaran sains pada jenjang SMA.

Menurut \textcite{hevner_design_2004}, DSR adalah paradigma pemecahan masalah
yang berupaya meningkatkan kapabilitas manusia dan organisasi dengan
menciptakan artefak baru, direpresentasikan dalam bentuk konstruk,
model, metode, dan instansiasi. Artefak yang dikembangkan dalam
penelitian ini meliputi:

\begin{enumerate}
  \item \textbf{Instansiasi (Purwarupa Sistem):} \textit{Knowledge Graph}
        kurikulum sains yang terinstansiasi pada basis data Neo4j beserta
        visualisasinya, sebagai purwarupa struktur data yang dapat ditelusuri
        dan diakses secara komputasional.
  \item \textbf{Metode (Algoritma dan \textit{Pipeline}):} Algoritma
        ekstraksi pengetahuan terstruktur menggunakan LLM dengan
        \textit{output} berbasis ontologi, serta \textit{pipeline} KG
        \textit{Construction} dan \textit{Completion} yang diuraikan pada
        Subbab~\ref{subsec:kerangka-konsep}.
\end{enumerate}

Secara metodologis, penelitian ini dikategorikan sebagai penelitian
terapan (\textit{applied research}) sebagaimana dirangkum pada
Tabel~\ref{tab:karakteristik}.

\begin{longtable}{@{}ll@{}}
  \caption{Karakteristik Metodologi Penelitian}
  \label{tab:karakteristik}\\
  \toprule
  \textbf{Karakteristik} & \textbf{Deskripsi} \\
  \midrule
  \endfirsthead
  \toprule
  \textbf{Karakteristik} & \textbf{Deskripsi} \\
  \midrule
  \endhead
  \midrule
  \endfoot
  \bottomrule
  \endlastfoot
  Jenis Penelitian & Terapan (\textit{Applied Research}) \\
  Pendekatan       & \textit{Design Science Research} (DSR) \\
  Metode           & Eksperimental dengan validasi kualitatif pakar \\
  Output           & Artefak sistem (Metode KGC dan Purwarupa Visualisasi) \\
  Paradigma        & Pragmatisme (berfokus pada utilitas dan pemecahan masalah) \\
\end{longtable}

Adopsi DSR dalam penelitian ini mengikuti kerangka kerja
\textcite{peffers_design_2007}, yang meliputi enam tahapan berikut.

\begin{enumerate}
  \item \textbf{Identifikasi Masalah:} Fragmentasi kurikulum sains
        lintas-disiplin (Fisika, Kimia, Biologi) menyebabkan siswa
        kesulitan melihat keterhubungan konsep.
  \item \textbf{Definisi Tujuan:} Mengembangkan sistem yang dapat
        menemukan relasi implisit antarkonsep lintas mata pelajaran.
  \item \textbf{Desain dan Pengembangan:} Membangun \textit{pipeline}
        ekstraksi--konstruksi--penyempurnaan \textit{knowledge graph}.
  \item \textbf{Demonstrasi:} Menerapkan artefak pada buku teks resmi
        Kemendikbudristek Fisika, Kimia, dan Biologi Kelas~XII.
  \item \textbf{Evaluasi:} Pengukuran metrik struktural (ADC,
\todo[inline,color=yellow!40]{[Soros] https://www.semantic-web-journal.net/system/files/swj3366.pdf}
        \textit{Modularity}), metrik ekstraksi, dan validasi pakar.
  \item \textbf{Komunikasi:} Dokumentasi dalam bentuk skripsi dan
        publikasi ilmiah.
\end{enumerate}

% ============================================================
\section{Lokasi dan Waktu Penelitian}
% ============================================================

\subsection{Lokasi Penelitian}

Penelitian ini dilaksanakan secara daring tanpa melibatkan studi
lapangan. Seluruh proses ekstraksi, konstruksi graf, dan pengumpulan
data validasi dilakukan melalui platform digital.

\subsection{Waktu Penelitian}

Penelitian berlangsung selama enam bulan, dari Januari hingga Juni 2026.
Rincian tahapan disajikan pada Tabel~\ref{tab:waktu}.

\begin{longtable}{@{}lll@{}}
  \caption{Rincian Waktu Penelitian}
  \label{tab:waktu}\\
  \toprule
  \textbf{Fase} & \textbf{Kegiatan} & \textbf{Periode} \\
  \midrule
  \endfirsthead
  \toprule
  \textbf{Fase} & \textbf{Kegiatan} & \textbf{Periode} \\
  \midrule
  \endhead
  \midrule
  \endfoot
  \bottomrule
  \endlastfoot
  1 & Studi literatur dan desain sistem
    & Januari -- Februari 2026 \\
  2 & Implementasi \textit{pipeline} ekstraksi
    & Februari -- Maret 2026 \\
  3 & Konstruksi \textit{Knowledge Graph}
    & \multirow{2}{*}{Maret -- April 2026} \\
  4 & Evaluasi dan validasi pakar \#1 & \\
  5 & \textit{Knowledge Graph Completion}
    & April 2026 \\
  6 & Evaluasi dan validasi pakar \#2
    & Mei 2026 \\
  7 & Analisis hasil dan penulisan
    & Mei -- Juni 2026 \\
\end{longtable}

% ============================================================
\section{Sumber Data dan Korpus Penelitian}
% ============================================================

Penelitian ini tidak menggunakan populasi manusia sebagai objek
eksperimen utama, melainkan dua jenis data yang saling melengkapi: data
\todo[inline,color=yellow!40]{[Tegar Wahyu] pindah ke 3.3 | balasan: }
tekstual sebagai korpus ekstraksi dan data penilaian pakar sebagai dasar
evaluasi model.

\subsection{Data Primer (Korpus Ekstraksi)}\label{subsec:data-primer}

Data primer yang digunakan dibatasi pada materi sains tingkat SMA
Fase~F (Kelas~XII). Pemilihan dokumen dilakukan secara
\textit{purposive sampling} berdasarkan relevansinya dengan Kurikulum
Merdeka. Korpus penelitian terdiri dari tiga buku teks siswa
(\textit{e-book}) terbitan Kemendikbudristek untuk mata pelajaran Fisika,
Kimia, dan Biologi Fase~F, yang diperoleh dari repositori resmi
\url{https://buku.kemdikbud.go.id}. Pembatasan pada buku teks resmi
memastikan bahwa entitas materi dan relasi dasar yang diekstrak sistem
benar-benar mencerminkan standar materi yang diajarkan secara nasional.

\subsection{Data Validasi Pakar}\label{subsec:data-validasi}
\todo[inline,color=yellow!40]{[Soros] Data B}

Data validasi diperoleh melalui penilaian kualitatif oleh panel pakar
terhadap \textit{triple} relasi---khususnya interkoneksi
lintas-disiplin---yang dihasilkan sistem KGC. Setiap \textit{triple}
dinilai menggunakan empat kategori penilaian yang telah didefinisikan pada
Subbab~\ref{subsec:validasi-manusia}, yaitu Benar (1,0), Sebagian Benar (0,5), Kurang Konteks
(0,25), dan Salah (0,0). Selain memberikan penilaian, pakar dapat
menambahkan komentar dan \textit{missing triples} yang dianggap esensial
namun tidak ditemukan sistem.

Pengumpulan data validasi dilakukan melalui aplikasi web KG Review App
yang dikembangkan khusus untuk penelitian ini (arsitektur teknis dibahas
pada Subbab~\ref{sec:instrumen-evaluasi}). Data penilaian akhir, yakni label
konsensus setelah resolusi ketidaksepakatan, menjadi
\textit{ground truth} untuk menghitung Fuzzy Precision, Recall, dan
F1-Score, serta menguji tingkat kesepakatan antarpenilai pada Fase~1
menggunakan Gwet's AC1.

% ============================================================
\section{Subjek Evaluasi dan Panel Pakar}
% ============================================================

Evaluasi performa sistem menggunakan pendekatan \textit{human
evaluation}. Subjek penelitian adalah panel pakar domain yang bertugas
sebagai \textit{reviewer} untuk memvalidasi akurasi ontologis dan
relevansi pedagogis relasi lintas-disiplin yang dihasilkan sistem.

Penelitian ini melibatkan enam \textit{reviewer} yang mewakili mata
pelajaran Fisika, Kimia, dan Biologi (dua \textit{reviewer} per mata
pelajaran). Jumlah dan kriteria pakar ditetapkan secara
\textit{purposive sampling} berdasarkan tiga pertimbangan metodologis
berikut.

\begin{enumerate}
  \item \textbf{Pengukuran Reliabilitas Multi-Rater:} Pada Fase~1, skema
        evaluasi mengharuskan satu \textit{triple} dalam-buku divalidasi
        oleh dua pakar independen, yang memungkinkan perhitungan
        \textit{inter-rater agreement} menggunakan Gwet's AC1. Adapun
        \textit{edge} lintas-mapel pada Fase~2 dinilai bersama melalui
        diskusi panel.

  \item \textbf{Beban Kognitif Penilaian:} Jumlah \textit{triple}
        relasi per mata pelajaran yang dihasilkan sistem melebihi 100
        item. Jumlah pakar dibatasi agar tidak menurunkan ketelitian
        akibat beban kognitif yang berlebihan selama proses anotasi.

  \item \textbf{Keketatan Kualifikasi:} \textit{Reviewer} harus
        merupakan mahasiswa aktif atau lulusan dengan latar belakang
        jurusan ilmu murni atau pendidikan yang linier (Biologi, Kimia,
        atau Fisika). Penggunaan mahasiswa dengan keahlian domain yang
        linier mengikuti preseden studi anotasi KG pendidikan yang
        mengutamakan kedalaman keilmuan spesifik dibanding pengalaman
        praktis umum \parencite{wei_kicgpt_2023}. Ketentuan ini secara alami
        membatasi ketersediaan populasi yang memenuhi kualifikasi.
\end{enumerate}

% ============================================================
\section{Kerangka Konsep}\label{subsec:kerangka-konsep}
% ============================================================

Kerangka konsep penelitian ini menggambarkan \textit{pipeline}
pemrosesan data dari dokumen kurikulum hingga \textit{Knowledge Graph
Completion}. Sistem dibangun dengan arsitektur empat tahap yang saling
terintegrasi untuk mengekstraksi, membangun, dan menyempurnakan jaringan
pengetahuan sains, serta dua fase evaluasi untuk memastikan validitasnya.

\subsection{Diagram Alur Penelitian}

\textit{Pipeline} penelitian terbagi ke dalam empat tahap eksekusi dan
dua fase evaluasi sebagaimana diilustrasikan pada
Gambar~\ref{fig:pipeline}.

\subfile{pipeline-kgc}

\noindent\textbf{Tahap-Tahap Eksekusi:}

\begin{enumerate}
  \item \textbf{Tahap 1 --- Pemrosesan Dokumen (\textit{Ingestion}):}
        Mengubah buku teks PDF menjadi format terstruktur, meliputi
        segmentasi hierarki bab, ekstraksi terminologi glosarium, dan
        pemotongan teks (\textit{sentence chunking}).

  \item \textbf{Tahap 2 --- Ekstraksi Pengetahuan (\textit{LLM
        Extraction}):} Menggunakan LLM terpandu
        (\textit{schema-constrained}) untuk mengekstrak entitas konsep
        dan memetakan relasi dalam-buku, divalidasi dengan terminologi
        glosarium.

  \item \textbf{Tahap 3 --- Konstruksi Graf (\textit{Graph
        Construction}):} Menginstansiasi hierarki struktural dan relasi
        semantik dalam-buku ke dalam basis data \textit{Knowledge Graph}
        (Neo4j).

  \item \textbf{Tahap 4 --- Penyempurnaan Graf (\textit{Graph
        Completion}):} Membangun vektor semantik (\textit{concept
        embedding}) dari konsep yang ada, mencari kemiripan lintas-buku
        menggunakan ANN, dan menggunakan LLM sebagai pengklasifikasi
        untuk memprediksi relasi antardisiplin.
\end{enumerate}

\noindent\textbf{Fase-Fase Evaluasi:}

\begin{enumerate}
  \item \textbf{Fase Evaluasi 1 --- Validasi Intra-Mapel \&
        \textit{LLM-as-a-judge}:} Diterapkan setelah Tahap~3. Kerangka
        kerja \textit{LLM-as-a-judge} (Gambar~\ref{fig:llm-judge})
        digunakan untuk mengevaluasi kumpulan \textit{triple} dalam satu
        mata pelajaran. Jika dua pakar tidak sepakat, AI Judge dipanggil
        sebagai penengah berbasis konteks buku teks sumber. Jika pakar
        mengidentifikasi \textit{triple} yang terlewat, LLM mengevaluasi
        kelayakan usulan penambahan sebelum diintegrasikan ke dalam KG.

  \item \textbf{Fase Evaluasi 2 --- Validasi Inter-Mapel \& Diskusi
        Pakar:} Diterapkan setelah Tahap~4. Memverifikasi \textit{edge}
        lintas-disiplin baru yang dihasilkan proses KGC. Mengingat
        sifat relasi yang implisit, konflik tidak diselesaikan oleh AI,
        melainkan melalui diskusi panel yang melibatkan seluruh pakar
        untuk mencapai konsensus keilmuan.
\end{enumerate}

\subfile{llm-judge}

\subsection{Skema Ontologi Knowledge Graph}
\label{subsec:skema-ontologi}

Ontologi dalam penelitian ini diposisikan sebagai model konseptual yang bersifat
\textit{deskriptif}, bukan \textit{preskriptif}: ontologi
merepresentasikan konsep-konsep sains beserta keterkaitan semantiknya
sebagaimana tersaji dalam buku teks Kurikulum Merdeka, tanpa
merekomendasikan urutan maupun jalur belajar tertentu. Relasi berarah
seperti \texttt{PRASYARAT\_UNTUK} karenanya dimaknai sebagai
ketergantungan konseptual yang terkandung dalam materi---bukan sebagai
anjuran kronologis bagi siswa.

Skema \textit{knowledge graph} diformalkan sebagai ontologi OWL untuk
memastikan konsistensi struktur dan vokabulari relasi antarstahap
\textit{pipeline} (ekstraksi, konstruksi, dan \textit{completion}).
Topologi keseluruhan---mencakup relasi struktural, pedagogis, dan
semantik---diilustrasikan pada Gambar~\ref{fig:ontologi}.

\subfile{onthology}

% ============================================================
\section{Parameter Sistem dan Definisi Operasional}
% ============================================================

\subsection{Parameter Sistem}
\label{subsec:parameter-sistem}

Parameter penelitian ini adalah variabel-variabel \textit{pipeline} yang
dapat dimodifikasi peneliti dan memengaruhi karakteristik \textit{knowledge
graph} yang dihasilkan. Variabel-variabel tersebut dirangkum pada
Tabel~\ref{tab:parameter}.

\begin{longtable}{@{}p{0.30\linewidth}p{0.35\linewidth}p{0.27\linewidth}@{}}
  \caption{Parameter Sistem \textit{Pipeline}}
  \label{tab:parameter}\\
  \toprule
  \textbf{Parameter} & \textbf{Deskripsi} & \textbf{Nilai/Level} \\
  \midrule
  \endfirsthead
  \toprule
  \textbf{Parameter} & \textbf{Deskripsi} & \textbf{Nilai/Level} \\
  \midrule
  \endhead
  \midrule
  \endfoot
  \bottomrule
  \endlastfoot
  Model LLM Ekstraksi \& Completion
    & Model bahasa besar untuk ekstraksi konsep/relasi dan
      penyempurnaan graf (KGC) lintas-disiplin
    & {\scriptsize \texttt{gemini-2.5-flash-lite} (akses Mei 2026)} \\
  Model LLM-as-a-judge
    & Model penengah ketidaksepakatan pakar pada \textit{pipeline} validasi
    & {\scriptsize \texttt{gemini-3.1-pro-preview} (akses Mei 2026)} \\
  Strategi \textit{Prompt Engineering}
    & Variasi instruksi sistem untuk ekstraksi
    & \textit{Schema-constrained few-shot} + \textit{context injection}
      + \textit{strict output constraints} \\
  Model \textit{Embedding}
    & Model representasi vektor untuk perhitungan similaritas
    & \texttt{gemini-embedding-001} (Gemini Embedding) \\
  Threshold Similaritas
    & Ambang batas \textit{cosine similarity} untuk KGC
    & Intra-buku: 0,80; Lintas-buku: 0,70 \\
  Parameter \textit{Chunking}
    & Pemotongan dokumen sebelum ekstraksi agar konteks stabil
    & \texttt{chunk\_size} = 800, \texttt{chunk\_overlap} = 200 \\
  Sumber Data
    & Domain materi yang diproses \textit{pipeline}
    & Biologi, Fisika, Kimia Kelas~XII \\
\end{longtable}

\subsection{Metrik Evaluasi}\label{subsec:metrik-evaluasi}

Metrik evaluasi dalam penelitian ini mencakup dimensi topologi graf dan
dimensi validasi pakar, sebagaimana dirangkum pada
Tabel~\ref{tab:metrik}. Definisi operasional lengkap setiap metrik
struktural telah dijabarkan pada Subbab~\ref{subsec:metrik-topologi}
dan~\ref{subsec:metrik-ekstraksi}; definisi
metrik validasi pakar pada Subbab~\ref{subsec:validasi-manusia}.

\begin{longtable}{@{}p{0.35\linewidth}p{0.37\linewidth}p{0.20\linewidth}@{}}
  \caption{Metrik Evaluasi Penelitian}
  \label{tab:metrik}\\
  \toprule
  \textbf{Variabel} & \textbf{Deskripsi} & \textbf{Satuan} \\
  \midrule
  \endfirsthead
  \toprule
  \textbf{Variabel} & \textbf{Deskripsi} & \textbf{Satuan} \\
  \midrule
  \endhead
  \midrule
  \endfoot
  \bottomrule
  \endlastfoot
  Jumlah Konsep Terekstrak
    & Volume konsep hasil ekstraksi per dokumen
    & Jumlah node \\
  Jumlah Relasi Terprediksi
    & Volume relasi lintas-buku hasil KGC
    & Jumlah edge \\
  \textit{Average Degree Centrality} (ADC)
    & Tingkat keterhubungan rata-rata; dibandingkan pre- dan post-KGC
    & Rata-rata derajat ($\geq 0$) \\
  \textit{Graph Density}
    & Kepadatan relasi graf; dibandingkan pre- dan post-KGC
    & Nilai 0--1 \\
  \textit{Modularity}
    & Keterpisahan komunitas antarmapel; dibandingkan pre- dan post-KGC
    & Nilai $-0{,}5$ hingga $1$ \\
  \textit{Description Completeness} (DC)
    & Kelengkapan deskripsi entitas
    & Nilai 0--1 \\
  \textit{Empty SubKonsep Rate} (ESR)
    & Proporsi entitas induk tanpa turunan
    & Nilai 0--1 \\
  \textit{Typed Relation Ratio} (TRR)
    & Proporsi relasi bertipe spesifik dari kosakata ontologi
    & Nilai 0--1 \\
  \textit{Fuzzy Precision}
    & Akurasi relasi berbasis skor kebenaran parsial (validasi pakar)
    & Nilai 0--1 \\
  \textit{Fuzzy Recall}
    & Cakupan relasi esensial berbasis skor kebenaran parsial (validasi pakar)
    & Nilai 0--1 \\
  F1-Score
    & Rata-rata harmonik \textit{Precision} dan \textit{Recall}
    & Nilai 0--1 \\
  Gwet's AC1
    & Kesepakatan antarpenilai yang dikoreksi terhadap kebetulan
    & Nilai $-1$ hingga $1$ \\
\end{longtable}

\subsection{Definisi Operasional}\label{subsec:definisi-operasional}

Tabel~\ref{tab:definisi} merangkum definisi operasional istilah-istilah
teknis utama yang digunakan sepanjang penelitian ini.

\begin{longtable}{@{}p{0.28\linewidth}p{0.65\linewidth}@{}}
  \caption{Definisi Operasional Istilah Penelitian}
  \label{tab:definisi}\\
  \toprule
  \textbf{Istilah} & \textbf{Definisi Operasional} \\
  \midrule
  \endfirsthead
  \toprule
  \textbf{Istilah} & \textbf{Definisi Operasional} \\
  \midrule
  \endhead
  \midrule
  \endfoot
  \bottomrule
  \endlastfoot
  \textit{Knowledge Graph}
    & Struktur data graf dalam Neo4j yang terdiri dari node berlabel
      \texttt{Grade/Subject}, \texttt{Chapter}, \texttt{Subtopic},
      \texttt{Concept}, dan \texttt{ConceptTarget}, serta \textit{edge}
      berupa relasi struktural (\texttt{HAS\_CHAPTER},
      \texttt{HAS\_CONCEPT}, dll.) maupun relasi semantik
      dalam-buku dan lintas-buku. \\
  Konsep (\texttt{Concept})
    & Entitas pengetahuan Fisika/Kimia/Biologi yang memiliki atribut
      \texttt{name} dan \texttt{description}. \\
  \texttt{ConceptTarget}
    & Node sementara yang menampung \textit{forward reference}---konsep
      yang disebut sebagai target relasi namun belum diekstrak secara
      lengkap; dipromosi menjadi \texttt{Concept} pada iterasi
      berikutnya. \\
  Relasi Struktural
    & Hubungan hierarkis dalam graf: \texttt{Grade/Subject} $\to$
      \texttt{Chapter} $\to$ \texttt{Subtopic} $\to$
      \texttt{Concept}. \\
  Relasi Implisit
    & Hubungan antarkonsep yang tidak tersurat dalam teks satu buku,
      diinferensikan melalui kemiripan semantik dan klasifikasi LLM. \\
  KGC (\textit{Knowledge Graph Completion})
    & Proses prediksi \textit{triple} (head, relation, tail) yang
      hilang menggunakan \textit{embedding similarity} dan klasifikasi
      LLM untuk menghasilkan relasi lintas-buku. \\
  \textit{Embedding}
    & Representasi vektor numerik dari deskripsi konsep yang dihasilkan
      model \textit{embedding} (\texttt{gemini-embedding-001}, dimensi:
      3072). \\
  \textit{Cosine Similarity}
    & Ukuran kemiripan antara dua vektor \textit{embedding}, dihitung
      sebagai \textit{dot product} dibagi hasil kali norma kedua vektor.\\
  ANN (\textit{Approximate Nearest Neighbor})
    & Algoritma pencarian tetangga terdekat berbasis HNSW pada Neo4j
      \textit{Vector Index} untuk mempercepat identifikasi kandidat
      relasi lintas-dokumen tanpa perbandingan brute-force
      $O(n^2)$. \\
  \textit{In-Context Learning} (ICL)
    & Strategi pemanfaatan LLM dengan menyertakan demonstrasi struktur
      ontologi di dalam instruksi \textit{prompt} tanpa memerlukan
      \textit{fine-tuning}. \\
\end{longtable}

% ============================================================
\section{Teknik Pengumpulan Data}
\todo[inline,color=yellow!40]{[Tegar Wahyu] pindahin ke detail tahapan | balasan: }
% ============================================================

Penelitian ini menggunakan dua jenis data yang telah diuraikan pada
Subbab~\ref{subsec:data-primer} dan~\ref{subsec:data-validasi}: data primer
\todo[inline,color=yellow!40]{[Soros] Data A}
kurikulum (dimanfaatkan pada Tahap~1--4) dan data validasi pakar (digunakan
pada kedua fase evaluasi). Karena sumber dan cakupan kedua jenis data tersebut
telah dibahas, bagian ini berfokus pada teknik pengumpulan data validasi pakar.

\subsection{Cara Pengumpulan Data}\label{subsec:cara-pengumpulan}
\todo[inline,color=yellow!40]{[Soros] cek}

Pengumpulan data validasi dilakukan melalui aplikasi web KG Review App
yang dikembangkan menggunakan Next.js, dipilih karena kemampuannya
menangani navigasi kuesioner berskala besar (lebih dari 100
\textit{triple} per mata pelajaran) yang tidak dapat diakomodasi oleh
Google Forms. Alur pengumpulan data adalah sebagai berikut.
\todo[inline,color=yellow!40]{[Tegar Wahyu] alur masukin detail tahap | balasan: }

\begin{enumerate}
  \item Peneliti mengunggah hasil \textit{knowledge graph} dalam format
        JSON ke panel admin aplikasi.
  \item Peneliti mempublikasikan mata pelajaran sehingga dapat diakses
        \textit{reviewer}.
  \item \textit{Reviewer} menerima kredensial \textit{login} dari
        peneliti.
  \item \textit{Reviewer} masuk ke aplikasi, memilih mata pelajaran yang
        ditugaskan, dan menilai setiap \textit{triple} menggunakan empat
        kategori penilaian (Subbab~\ref{subsec:validasi-manusia}).
  \item \textit{Reviewer} dapat menambahkan \textit{missing triples} dan
        memberikan komentar per relasi.
  \item Setelah seluruh \textit{triple} dinilai, \textit{reviewer}
        mengonfirmasi penyelesaian \textit{review}.
\end{enumerate}

% ============================================================
\section{Pengolahan Data}\label{sec:pengolahan-data}
\todo[inline,color=yellow!40]{[Tegar Wahyu] udah ada di metrik | balasan: }
% ============================================================

Setelah seluruh \textit{reviewer} menyelesaikan penilaian, data rating
diolah melalui tahapan berikut.

\begin{enumerate}
  \item \textbf{Ekspor dan Agregasi:} Data penilaian diekspor dari Redis
        dalam format JSON per mata pelajaran. Pada Fase~1, setiap
        \textit{triple} dalam-buku memiliki dua rating independen yang
        diagregasi menjadi satu label konsensus berdasarkan kerangka
        evaluasi \textit{hybrid}. Pada Fase~2, setiap \textit{edge}
        \textit{completion} lintas-mapel dinilai melalui diskusi panel
        sehingga langsung menghasilkan satu label konsensus tunggal.

  \item \textbf{Penetapan \textit{Ground Truth}:} Jika dua
        \textit{reviewer} sepakat, label konsensus digunakan langsung.
        Jika tidak sepakat, AI Judge (LLM berbasis konteks buku teks)
        menentukan label akhir---kecuali pada Fase~2, yang
        menggunakan diskusi panel pakar secara keseluruhan.

  \item \textbf{Perhitungan Metrik Pakar:} Fuzzy Precision, Recall, dan
        F1-Score dihitung per bab per mata pelajaran menggunakan
        formulasi pada Subbab~\ref{subsec:validasi-manusia}, dengan label konsensus sebagai
        \textit{ground truth}.

  \item \textbf{Perhitungan Kesepakatan:} Gwet's AC1 dihitung khusus pada
        Fase~1, mengukur konsistensi \textit{inter-rater} sebelum
        konsolidasi menggunakan label mentah dari masing-masing
        \textit{reviewer}. \textit{Edge} \textit{completion} pada Fase~2
        dinilai melalui diskusi panel sehingga hanya menghasilkan satu
        label konsensus dan tidak masuk dalam perhitungan
        \textit{inter-rater agreement}.

  \item \textbf{Kalkulasi Metrik Struktural:} Metrik topologi graf
        (ADC, \textit{Density}, \textit{Modularity}, TRR, DC) dihitung
        melalui kueri Cypher deterministik di Neo4j Aura pada kondisi
        pre-KGC dan post-KGC.
\end{enumerate}

% ============================================================
\section{Analisis Data}
\todo[inline,color=yellow!40]{[Soros] Cek}
% ============================================================

Analisis data dilakukan melalui tiga pendekatan yang saling melengkapi,
dirancang untuk menjawab rumusan masalah penelitian secara komprehensif
dari sisi kuantitatif objektif maupun validitas pedagogis.

\subsection{Analisis Struktural}

Analisis struktural menafsirkan perubahan karakteristik topologi
\textit{Knowledge Graph} untuk mengevaluasi dampak proses KGC terhadap
kualitas interkoneksi antarkonsep. Berdasarkan metrik struktural yang
dihitung pada tahap pengolahan data (Subbab~\ref{sec:pengolahan-data}),
hasil perbandingan kondisi pre-KGC dan post-KGC disajikan dalam bentuk
tabel berpasangan per mata pelajaran. Peningkatan ADC dan \textit{Density}
setelah KGC diinterpretasikan sebagai indikator bahwa relasi implisit
yang diprediksi LLM berhasil memperkaya kepadatan jaringan pengetahuan.
Sebaliknya, penurunan \textit{Modularity} yang terlalu drastis dapat
mengindikasikan bahwa batas komunitas antarbab menjadi kabur akibat
relasi lintas-disiplin yang berlebihan, sehingga perlu dikaji lebih
lanjut melalui validasi pakar.

\subsection{Analisis Validasi Pakar}

Analisis ini menafsirkan akurasi dan kelengkapan relasi yang dihasilkan KG
dari metrik validasi pakar yang dihitung pada Subbab~\ref{sec:pengolahan-data}.
Nilai Gwet's AC1 dibaca sebagai keandalan kesepakatan antarpakar pada Fase~1,
sedangkan Fuzzy Precision, Recall, dan F1-Score \parencite{harju_evaluating_2023}
ditafsirkan sebagai kinerja sistem terhadap \textit{ground truth} konsensus
(Subbab~\ref{subsec:validasi-manusia}).

\subsection{Analisis Komparatif}

Analisis komparatif membandingkan kualitas \textit{Knowledge Graph} di
antara ketiga mata pelajaran untuk mengidentifikasi perbedaan
karakteristik struktur pengetahuan antardisiplin sains dalam konteks
Kurikulum Merdeka.

Perbandingan dilakukan pada dua dimensi. Pertama, dimensi struktural
membandingkan nilai metrik graf (ADC, \textit{Density},
\textit{Modularity}, TRR) antarmata pelajaran pada kondisi post-KGC.
Kedua, dimensi validasi membandingkan nilai Precision, Recall, dan F1
antarmata pelajaran untuk mengidentifikasi disiplin mana yang paling
diuntungkan oleh pendekatan LLM-Assisted KGC.

Triangulasi antara metrik struktural dan hasil validasi pakar dilakukan
untuk menguji koherensi keduanya: apabila mata pelajaran dengan ADC
tinggi juga memperoleh nilai Precision tinggi dari \textit{reviewer},
maka ini memperkuat argumen bahwa kepadatan koneksi yang dihasilkan
sistem memang mencerminkan interkoneksi yang valid secara pedagogis,
bukan sekadar artefak komputasional. Hasil analisis disajikan dalam
tabel ringkasan tiga kolom (Fisika, Kimia, Biologi) yang memuat seluruh
metrik secara berdampingan.

\ifSubfilesClassLoaded{\printbibliography[title={Daftar Pustaka}]}{}

\end{document}
